\section{Ethics}\label{sec:ethics}

The experiment touches only public records: published opinions,
docket metadata, and audio of public oral arguments, all released
by the Court and mirrored by public-interest archives. No private
data, no sealed-at-source material, and no personal data beyond
the public professional conduct of public officials is used. The
justices are public figures acting in an official capacity; the
measured quantities (vote directions, agreement rates, writing
style of signed opinions) are matters of public record and public
commentary.

Two commitments go beyond data source. First, counterfactual
labels are labeled as counterfactuals: when the site or the models
display ``what the baseline would have voted'' for a case a
justice never heard (bench turnover makes this common), the row is
marked as a simulated call, never presented as a prediction the
justice could have made. Second, the project is explicitly
non-commercial: no paywall, no product, no data licensing ---
the corpus, the scripts, the trained adapters, and this paper are
published for replication and criticism. The zero-budget
constraint is partly ethical in origin: results that can only be
reproduced with institutional compute are claims that can only be
audited by institutions, and judicial-prediction claims should be
auditable by anyone.

A final note on downstream use. A high-accuracy vote predictor
would be commercially interesting to lobbyists and litigators; the
current state of the project---a 63.7\% ideological bar that
structured models fail to beat---is, we note dryly, a result
with no immediate commercial application, and we are in no hurry
to change that.
